\chapter{Foundations}\label{record-021}

\noindent{\footnotesize\color{BookMuted}\path{docs/design-system/foundations.md}\par}\smallskip

This document defines the shared design language behind Bijou.

It is not an API document. It is the policy layer that should shape API design, examples, and component usage.

\section{1. Color}\label{record-021-color}

\subsection{Purpose}\label{record-021-purpose}

Color in Bijou should communicate:

\begin{itemize}
\tightlist
\item
  status
\item
  hierarchy
\item
  separation
\item
  focus
\item
  interruption level
\end{itemize}

Color should not carry meaning alone. Every important state must still survive \texttt{n\allowbreak{}o\allowbreak{}C\allowbreak{}o\allowbreak{}l\allowbreak{}o\allowbreak{}r}, \texttt{p\allowbreak{}i\allowbreak{}p\allowbreak{}e}, and \texttt{a\allowbreak{}c\allowbreak{}c\allowbreak{}e\allowbreak{}s\allowbreak{}s\allowbreak{}i\allowbreak{}b\allowbreak{}l\allowbreak{}e} runs.

\subsection{Color roles}\label{record-021-color-roles}

Bijou should treat tokens as belonging to a small set of roles:

\begin{itemize}
\tightlist
\item
  \textbf{foreground}

  \begin{itemize}
  \tightlist
  \item
    primary text, secondary text, subdued text
  \end{itemize}
\item
  \textbf{surface}

  \begin{itemize}
  \tightlist
  \item
    canvas, elevated region, overlay, inset region
  \end{itemize}
\item
  \textbf{status}

  \begin{itemize}
  \tightlist
  \item
    success, warning, error, info, neutral
  \end{itemize}
\item
  \textbf{accent}

  \begin{itemize}
  \tightlist
  \item
    active/focused selection, primary affordance
  \end{itemize}
\item
  \textbf{border/separator}

  \begin{itemize}
  \tightlist
  \item
    grouping and structural division
  \end{itemize}
\end{itemize}

The exact API-level mapping from these roles to Bijou's shipped token families is defined in \href{https://github.com/flyingrobots/bijou/blob/cbd70612140e7a450d764b3c7be1c6af878b99bc/docs/design-system/theme-tokens.md}{Theme Token Vocabulary}.

\subsection{Color rules}\label{record-021-color-rules}

\begin{itemize}
\tightlist
\item
  Use status color to explain state, not to decorate.
\item
  Use accent color for focus and active affordances, not for arbitrary emphasis.
\item
  Use background fills sparingly and deliberately. A background should indicate containment, elevation, or interruption.
\item
  Avoid using more than one strong status tone in the same local region unless comparing states is the point.
\item
  In \texttt{n\allowbreak{}o\allowbreak{}C\allowbreak{}o\allowbreak{}l\allowbreak{}o\allowbreak{}r}, every state must still read correctly from words, symbols, order, or borders.
\end{itemize}

\subsection{Terminal-specific guidance}\label{record-021-terminal-specific-guidance}

\begin{itemize}
\tightlist
\item
  Strong color contrast is expensive in a dense terminal. Use saturated status tones for small focal areas, not for entire screens.
\item
  Large filled regions should be reserved for overlays, selected/focused states, or major grouping moves.
\item
  Bright foreground on dark background is the default mental model in many terminals, but docs should not assume dark mode semantics.
\end{itemize}

\section{2. Spacing and rhythm}\label{record-021-spacing-and-rhythm}

Carbon uses a 2x grid. Bijou needs a terminal-first equivalent.

\subsection{Bijou spacing model}\label{record-021-bijou-spacing-model}

\begin{itemize}
\tightlist
\item
  The atomic unit is \textbf{1 cell}.
\item
  The default rhythm unit is \textbf{2 cells}.
\item
  Dense internal padding can use 1-cell increments.
\item
  Structural gaps between sections, panes, and overlay edges should prefer the 2-cell rhythm.
\end{itemize}

\subsection{Default guidance}\label{record-021-default-guidance}

\begin{itemize}
\tightlist
\item
  inline gap: \texttt{1}
\item
  local group gap: \texttt{1} or \texttt{2}
\item
  section gap: \texttt{2}
\item
  pane or card padding: usually \texttt{1} vertically, \texttt{2} horizontally
\item
  modal margin from viewport edge: minimum \texttt{1}, preferred \texttt{2}
\item
  stacked transient gap: \texttt{1}
\end{itemize}

\subsection{Why this matters}\label{record-021-why-this-matters}

Without a declared rhythm, terminal UIs become visually noisy very fast. Bijou should feel intentionally packed, not cramped.

\section{3. Text hierarchy and emphasis}\label{record-021-text-hierarchy-and-emphasis}

Terminal typography is mostly:

\begin{itemize}
\tightlist
\item
  placement
\item
  spacing
\item
  boldness
\item
  token choice
\item
  borders/dividers
\item
  line length
\end{itemize}

not font families or font sizes.

\subsection{Hierarchy levels}\label{record-021-hierarchy-levels}

\begin{itemize}
\tightlist
\item
  \textbf{page/shell title}

  \begin{itemize}
  \tightlist
  \item
    strongest label in the local frame
  \end{itemize}
\item
  \textbf{section title}

  \begin{itemize}
  \tightlist
  \item
    container or region label
  \end{itemize}
\item
  \textbf{body}

  \begin{itemize}
  \tightlist
  \item
    default content
  \end{itemize}
\item
  \textbf{supporting text}

  \begin{itemize}
  \tightlist
  \item
    help text, metadata, timestamps, hints
  \end{itemize}
\end{itemize}

\subsection{Emphasis rules}\label{record-021-emphasis-rules}

\begin{itemize}
\tightlist
\item
  Prefer structure before decoration.

  \begin{itemize}
  \tightlist
  \item
    use spacing, grouping, and labeling first
  \end{itemize}
\item
  Bold should mark titles, active labels, or short key facts.
\item
  Dim/subdued text should carry metadata, not primary instructions.
\item
  Avoid combining too many modifiers at once.

  \begin{itemize}
  \tightlist
  \item
    bold + color + underline + inverse is usually design panic, not design clarity
  \end{itemize}
\end{itemize}

\subsection{Content density}\label{record-021-content-density}

\begin{itemize}
\tightlist
\item
  Tight density is acceptable in terminals, but meaning must still scan.
\item
  Long prose should wrap and breathe.
\item
  Dense data views should prefer alignment and scanability over decorative framing.
\end{itemize}

\section{4. Overflow}\label{record-021-overflow}

Bijou should default to \textbf{wrap for prose-like content} and \textbf{truncate or clip only where layout semantics demand it}.

\subsection{Default policy}\label{record-021-default-policy}

\begin{itemize}
\tightlist
\item
  prose, supporting text, and descriptions: wrap
\item
  labels inside constrained shell chrome: clip or truncate
\item
  tabular cells: wrap when the row can grow, truncate when preserving scanability is more important
\item
  activity streams and logs: do not silently destroy ordering to preserve prettiness
\end{itemize}

\subsection{Why}\label{record-021-why}

Truncation is often a hidden data-loss bug disguised as layout discipline.

\subsection{Required controls}\label{record-021-required-controls}

Families that render text in constrained spaces should expose overflow choices where practical:

\begin{itemize}
\tightlist
\item
  \texttt{w\allowbreak{}r\allowbreak{}a\allowbreak{}p}
\item
  \texttt{t\allowbreak{}r\allowbreak{}u\allowbreak{}n\allowbreak{}c\allowbreak{}a\allowbreak{}t\allowbreak{}e}
\item
  \texttt{c\allowbreak{}l\allowbreak{}i\allowbreak{}p}
\end{itemize}

Not every component needs every mode, but the library should speak one consistent overflow language.

\section{5. Motion}\label{record-021-motion}

Motion should clarify:

\begin{itemize}
\tightlist
\item
  entry
\item
  exit
\item
  relationship between before and after
\item
  focus shifts
\item
  interruption level
\end{itemize}

Motion should not exist just because the runtime can animate.

\subsection{Motion rules}\label{record-021-motion-rules}

\begin{itemize}
\tightlist
\item
  Transient surfaces may animate in and out.
\item
  Page or pane transitions should reinforce spatial or logical movement.
\item
  Fast repetitive interactions should avoid theatrical motion.
\item
  Any animation must have a deterministic non-animated truth underneath it.
\end{itemize}

\subsection{Practical policy}\label{record-021-practical-policy}

\begin{itemize}
\tightlist
\item
  notifications: short edge-based slide or fade motion
\item
  shell/page transitions: modest, directional, and optional
\item
  loading states: subtle looping motion
\item
  structural resize: prefer stable layout over decorative flourish
\end{itemize}

\section{6. Elevation and interruption}\label{record-021-elevation-and-interruption}

Terminals do not have shadows in the usual sense. Elevation must be conveyed through:

\begin{itemize}
\tightlist
\item
  border treatment
\item
  fill treatment
\item
  separation from background
\item
  reserved space
\item
  z-order and animation
\end{itemize}

\subsection{Interruption ladder}\label{record-021-interruption-ladder}

From least interruptive to most interruptive:

\begin{enumerate}
\def\labelenumi{\arabic{enumi}.}
\tightlist
\item
  inline help or note
\item
  badge or subtle status
\item
  alert
\item
  toast / notification
\item
  drawer
\item
  modal
\end{enumerate}

This ladder should show up consistently in docs and examples.

\section{7. Focus and active state}\label{record-021-focus-and-active-state}

Bijou is keyboard-first.

The user should always be able to tell:

\begin{itemize}
\tightlist
\item
  what is focused
\item
  what is active
\item
  what is selected
\item
  what is merely present
\end{itemize}

\subsection{Focus cues}\label{record-021-focus-cues}

Use some combination of:

\begin{itemize}
\tightlist
\item
  accent color
\item
  border or gutter treatment
\item
  active label treatment
\item
  position within a structured list
\end{itemize}

\subsection{Rules}\label{record-021-rules}

\begin{itemize}
\tightlist
\item
  Do not rely on color alone for focus.
\item
  Focus cues should be local and unmistakable.
\item
  Shell-level focus and component-level focus should not compete visually at the same strength.
\end{itemize}

\section{8. Input modality policy}\label{record-021-input-modality-policy}

\subsection{Keyboard first}\label{record-021-keyboard-first}

Bijou should assume keyboard interaction is the baseline.

Every important interactive flow should work without a mouse.

\subsection{Mouse as enhancement}\label{record-021-mouse-as-enhancement}

Mouse input is valuable, but it should be:

\begin{itemize}
\tightlist
\item
  additive
\item
  discoverable
\item
  consistent
\item
  never the only path
\end{itemize}

\subsection{Pointer policy}\label{record-021-pointer-policy}

\begin{itemize}
\tightlist
\item
  Click should usually mirror an existing keyboard action.
\item
  Hover-only meaning should be minimal.
\item
  Drag should be reserved for clearly spatial interactions, such as split dividers.
\item
  Overlays should consume mouse events before background content.
\end{itemize}

\section{9. Package ownership}\label{record-021-package-ownership}

\subsection{\texorpdfstring{Core \texttt{@\allowbreak{}f\allowbreak{}l\allowbreak{}y\allowbreak{}i\allowbreak{}n\allowbreak{}g\allowbreak{}r\allowbreak{}o\allowbreak{}b\allowbreak{}o\allowbreak{}t\allowbreak{}s\allowbreak{}/\allowbreak{}b\allowbreak{}i\allowbreak{}j\allowbreak{}o\allowbreak{}u}}{Core }}\label{record-021-core-flyingrobotsbijou}

Use core for semantics that survive graceful degradation:

\begin{itemize}
\tightlist
\item
  status and structural components
\item
  forms and prompts
\item
  prose and log output
\item
  data views that still make sense as text
\end{itemize}

\subsection{\texorpdfstring{\texttt{@\allowbreak{}f\allowbreak{}l\allowbreak{}y\allowbreak{}i\allowbreak{}n\allowbreak{}g\allowbreak{}r\allowbreak{}o\allowbreak{}b\allowbreak{}o\allowbreak{}t\allowbreak{}s\allowbreak{}/\allowbreak{}b\allowbreak{}i\allowbreak{}j\allowbreak{}o\allowbreak{}u\allowbreak{}-\allowbreak{}t\allowbreak{}u\allowbreak{}i}}{}}\label{record-021-flyingrobotsbijou-tui}

Use TUI for:

\begin{itemize}
\tightlist
\item
  shells
\item
  overlays
\item
  scrolling
\item
  focus-heavy inspection
\item
  animation
\item
  notification systems
\item
  pointer-heavy or lifecycle-heavy behavior
\end{itemize}

\section{10. What this means for future library work}\label{record-021-what-this-means-for-future-library-work}

New component work should start by answering:

\begin{enumerate}
\def\labelenumi{\arabic{enumi}.}
\tightlist
\item
  Which pattern is this serving?
\item
  Is it semantic, interaction-layer, or render-path variation?
\item
  Does it belong in core or TUI?
\item
  What is the default overflow behavior?
\item
  How does it lower through non-interactive modes?
\item
  What is the Carbon-style ``when to use / when not to use'' story?
\end{enumerate}

